% 368_p_Werte_Lepton_Massen_En_ch.tex
% Johann Pascher, 19 September 2026
% From p0,p1,p2 to lepton masses: phi-skeleton and Galois correction factors

\providecommand{\BM}{\textbf{[B]}}
\providecommand{\KM}{\textbf{[K]}}
\providecommand{\SM}{\textbf{[S]}}
\providecommand{\EM}{\textbf{[E]}}

\chapter*{From $p_0,p_1,p_2$ to Lepton Masses: $\phi$-Skeleton and Galois Correction Factors}
\addcontentsline{toc}{chapter}{From $p_0,p_1,p_2$ to Lepton Masses}

\begin{abstract}
Doc.~293 provides the icosahedral redistribution weights $p_0=2/9$,
$p_1=(2{+}3\phi)/9$, $p_2=(5{-}3\phi)/9$. Doc.~367 shows that $p_0=\theta$
\emph{and} rewrites the Koide formula entirely in matrix elements of $R_5$ --
the mass bridge via $\theta$ is thereby obtained directly from the $A_5$ geometry.

This document investigates a second, purely algebraic route: the \emph{ratios}
$p_1/p_2$ and $p_0/p_2$ carry the $\phi$-skeleton of the lepton masses directly,
without a cosine function.
Main result: $p_1/p_2 = \phi^8$ and $p_0/p_2 = 2\phi^4$ are exact algebraic
identities \BM. The mass ratios $m_\tau/m_e$ and $m_\mu/m_e$ follow approximately
by multiplication with the Galois factors $74$ and $15$ from Doc.~338
(accuracy $0{.}03\,\%$ and $0{.}56\,\%$) \KM. The derivation of these Galois
factors from the $\mathbb{Z}_3$ framework remains open \SM.
\end{abstract}

\section*{Notation}
\addcontentsline{toc}{section}{Notation}

\textbf{Epistemic markers:}
\begin{itemize}[leftmargin=2em,itemsep=0pt,topsep=2pt]
  \item[\textbf{[B]}] algebraically proved, exact arithmetic
  \item[\textbf{[K]}] numerically confirmed (approximation, accuracy stated)
  \item[\textbf{[S]}] structurally observed, derivation open
  \item[\textbf{[E]}] established external mathematics
\end{itemize}

\medskip
\begin{tabular}{@{}lp{11cm}@{}}
\toprule
\textbf{Symbol} & \textbf{Meaning} \\
\midrule
$\phi$ & golden ratio $(1+\sqrt5)/2$ \\
$p_0,p_1,p_2$ & icosahedral redistribution weights from Doc.~293 \\
$\theta$ & Koide phase; $\theta=p_0=2/9$ (Doc.~367) \\
$a_k$ & Koide amplitude: $\sqrt{m_k}=M\,(1+\sqrt2\cos(\theta+2\pi k/3))$ \\
\bottomrule
\end{tabular}

\section{Starting Point}

Doc.~293 computes the redistribution weights of the electron mode $v_e=v_1$
under the icosahedral five-fold rotation $R_5\in A_5$:
\[
  p_0 = \frac{2}{9}, \qquad
  p_1 = \frac{2+3\phi}{9}, \qquad
  p_2 = \frac{5-3\phi}{9}, \qquad
  p_0+p_1+p_2=1.
\]
Doc.~367 establishes that $\theta=p_0=2/9$: the same quotient appears once
as Koide phase and once as transition probability. Doc.~367 has shown that the mass bridge via $\theta = p_0$ can be obtained
directly from the $A_5$ geometry. The question of this document is a different
one: do the \emph{ratios} $p_1/p_2$ and $p_0/p_2$ carry the $\phi$-skeleton
of the masses algebraically -- without cosine function, without $\theta$ as a
running variable?

\section{Exact Algebraic Identities}

\subsection{$p_1/p_2 = \phi^8$}

\begin{equation}\label{eq:p1p2}
  \frac{p_1}{p_2} = \frac{2+3\phi}{5-3\phi} = \frac{(7+3\sqrt5)/2}{(7-3\sqrt5)/2}
  = \frac{(7+3\sqrt5)^2}{49-45} = \frac{94+42\sqrt5}{4} = \frac{47+21\sqrt5}{2}.
\end{equation}
On the other hand: $\phi^4 = 3\phi+2$ (from $\phi^2=\phi+1$), so
$\phi^8 = (3\phi+2)^2 = 21\phi+13 = \tfrac{47+21\sqrt5}{2}$. Hence \BM:
\[
  \boxed{p_1/p_2 = \phi^8.}
\]

\subsection{$p_0/p_2 = 2\phi^4$}

\[
  \frac{p_0}{p_2} = \frac{2}{5-3\phi} = \frac{2(7+3\sqrt5)}{4} = \frac{7+3\sqrt5}{2}.
\]
And $2\phi^4 = 6\phi+4 = 7+3\sqrt5$. Hence \BM:
\[
  \boxed{p_0/p_2 = 2\phi^4.}
\]

\subsection{$p_1/p_0 = \phi^4/2$}

Direct consequence: $p_1/p_0 = \phi^8/(2\phi^4) = \phi^4/2$ \BM.
All three ratios carry only $\phi$-powers with rational prefactors.

\subsection{Amplitude Structure}

The square roots $3\sqrt{p_j}$ (amplitudes up to normalisation factor 3)
take exactly the icosahedral values \BM:
\[
  3\sqrt{p_0}=\sqrt2, \qquad 3\sqrt{p_1}=\phi^2, \qquad 3\sqrt{p_2}=\tfrac{1}{\phi^2}.
\]

\section{Approximate Connection to Mass Ratios}

\subsection{Precision limit: $\cos(2/9)$ is transcendental}

The exact Koide mass ratios at $\theta=2/9$ are:
\[
  \frac{m_\tau}{m_e} = \left(\frac{a_0}{a_1}\right)^2, \qquad
  \frac{m_\mu}{m_e} = \left(\frac{a_2}{a_1}\right)^2,
\]
with $a_k = 1+\sqrt2\cos(2/9+2\pi k/3)$. Since $2/9\notin\pi\mathbb{Q}$,
$\cos(2/9)$ is transcendental \EM. Consequently $m_\tau/m_e$ and $m_\mu/m_e$
cannot be expressed exactly as rational multiples of $\phi$-powers. All
connections below are approximate identities with stated accuracy.

\subsection{$m_\tau/m_e \approx 74\times\phi^8 = 74\times p_1/p_2$}

Numerically (with $\theta=2/9$ exact in the Koide formula) \KM:
\begin{align}
  \frac{m_\tau}{m_e}\bigl|_{\theta=2/9} &= 3477{.}47,\\
  74\times\phi^8 &= 3476{.}42,\\
  \text{deviation} &= 0{.}030\,\%.
\end{align}
The Galois factor $74=2\times37$, where $37\mid|GF(3^{18})^*|$, appears
in Doc.~338 as part of the correction factor $K_\text{frak}=74/75$ of the
$\alpha$ derivation. Its appearance here is a structural observation \SM;
the derivation from the $\mathbb{Z}_3$ framework is outstanding.

\subsection{$m_\mu/m_e \approx 30\times\phi^4 = 15\times p_0/p_2$}

\begin{align}
  \frac{m_\mu}{m_e}\bigl|_{\theta=2/9} &= 206{.}770,\\
  30\times\phi^4 &= 205{.}623,\\
  \text{deviation} &= 0{.}56\,\%.
\end{align}
Here the approximation is considerably less precise. The factor $15=3\times5$,
with $3=|GF(3)|$ and $5=\text{min.prim}(|GF(81)^*|)$, comes from the Galois
hierarchy of Doc.~338 \SM.

\section{Numerical Confirmation}

The check script \texttt{pruef\_368\_p\_massen.py} (2/python/Dok368\_Skripte)
confirms \KM:
\begin{enumerate}
  \item $p_1/p_2=\phi^8$ exactly (dev.\ $<10^{-35}$, 40 digits).
  \item $p_0/p_2=2\phi^4$ exactly.
  \item $p_1/p_0=\phi^4/2$ exactly.
  \item $3\sqrt{p_j}\in\{\sqrt2,\phi^2,1/\phi^2\}$ exactly.
  \item $74\times\phi^8$ matches $m_\tau/m_e|_{\theta=2/9}$ to $0{.}030\,\%$.
  \item $30\times\phi^4$ matches $m_\mu/m_e|_{\theta=2/9}$ to $0{.}56\,\%$.
  \item All three ratios carry only $\phi$-powers (no $\pi$, no $\cos$).
  \item $\varepsilon \approx 27\xi/12$ (0.49\,\% match).
  \item $74\times\phi^8\times(1+27\xi/12)$ matches $m_\tau/m_e$ to $0{.}00015\,\%$.
\end{enumerate}

\section{Galois Structure of $\varepsilon$: Observation and Limit}

\subsection{Numerical observation}

The deviation $\varepsilon$ from Section~3 lies close to a Galois-native
$\xi$-expression:
\begin{equation}\label{eq:eps_galois}
  \varepsilon \approx \frac{27\xi}{12}
  = \frac{|GF(27)|\cdot\xi}{n_{\phi,\mu}\cdot n_{\theta,e}}
  = \frac{3}{10000}, \qquad \text{dev.\ } 0{.}49\,\% \quad \KM
\end{equation}
All three factors -- $27 = |GF(27)|$, $\xi = 4/30000$ and
$12 = n_{\phi,\mu}\times n_{\theta,e} = 4\times3$ -- are already known
with status \BM\ from Doc.~338. Substituting:
\begin{equation}
  \frac{m_\tau}{m_e} \approx 74\times\phi^8\times\!\left(1+\frac{27\xi}{12}\right)
  \quad (0{.}00015\,\%) \quad \KM
\end{equation}

\subsection{Limit of the analytical proof path}

A proof that $\varepsilon = 27\xi/12$ holds \emph{exactly} cannot be
established via the FFGFT formula. The test attempt:

From Doc.~352: $m_\tau/m_e(\text{bare}) = (25/12)\cdot\xi^{-5/6}$
(from winding numbers $r_\tau/r_e = 25/12$, $p_\tau-p_e=-5/6$), with
fraktal correction coefficient $117{.}4\cdot\xi$. A linearisation
$(25/12)\cdot\xi^{-5/6} \approx 74\times\phi^8\times(1+c_\text{bare}\cdot\xi)$
yields $c_\text{bare} \approx 119{.}1$.

The net coefficient is:
\[
  c_\text{bare} - 117{.}4 \approx 1{.}72 \neq \frac{27}{12} = 2{.}25
  \quad (23\,\%\text{ discrepancy})
\]
The approach fails for two reasons:
\begin{enumerate}
  \item $\xi^{-5/6}$ is not a polynomial in $\xi$: the expansion of
    $(25/12)\cdot\xi^{-5/6}$ around the fixed point $74\times\phi^8$ yields
    a power series in $\xi^{1/6}$, not in $\xi$.
  \item The coefficient $117{.}4$ in Doc.~352 is itself numerical \KM\
    (rounded value); both sides carry uncertainties that add up to $\sim23\,\%$.
\end{enumerate}

\paragraph{Status.}
The identification $\varepsilon \approx 27\xi/12$ is a numerical
observation \KM\ (0.49\,\% deviation). All factors are Galois-native from
Doc.~338. An exact analytical derivation -- showing why the net coefficient
is precisely $27/12$ -- requires a perturbation expansion in $\xi^{1/6}$
that goes beyond the current state of the corpus. This step remains \SM.

\section{Summary: What Is Computable from the $p$-Values}

\begin{center}
\begin{tabular}{@{}p{0.42\linewidth}p{0.12\linewidth}p{0.38\linewidth}@{}}
\toprule
\textbf{Statement} & \textbf{Status} & \textbf{Source} \\
\midrule
$p_1/p_2 = \phi^8$, $\;p_0/p_2 = 2\phi^4$ & \BM & $A_5$ geometry, Doc.~293 \\
$3\sqrt{p_j}\in\{\sqrt2,\phi^2,1/\phi^2\}$ & \BM & $A_5$ geometry, Doc.~293 \\
$m_\tau/m_e \approx 74\times p_1/p_2$ & \KM & Galois factor 74, Doc.~338 \\
$m_\mu/m_e \approx 30\times p_0/p_2$ & \KM & Galois factor 15, Doc.~338 \\
$\varepsilon \approx 27\xi/12$ & \KM & Galois$+\xi$, Doc.~338 \\
Derivation of factors 74, 30 from $p$-values & \SM & open \\
Analytic proof $\varepsilon = 27\xi/12$ exactly & \SM & open \\
\bottomrule
\end{tabular}
\end{center}

\paragraph{What follows directly from the $p$-values.}
The $p$-values $p_0,p_1,p_2$ from the $A_5$ geometry (Doc.~293)
deliver exactly the \emph{$\phi$-skeleton} of the lepton masses: the
ratios $p_1/p_2 = \phi^8$ and $p_0/p_2 = 2\phi^4$ are purely geometric
identities with no free parameter.

\paragraph{What requires additional input.}
Converting to numerical mass ratios requires the Galois factors 74 and 30.
These come from the $GF(3^n)$ structure of Doc.~338, not from the $p$-values
themselves. The direct route
\[
  p_0,p_1,p_2 \;\xrightarrow{\text{exact}}\; \phi\text{-skeleton}
  \;\xrightarrow{\text{Galois, Doc.\,338}}\; m_\tau/m_e,\;m_\mu/m_e
\]
is fully described, but the second step is observation, not derivation.
A fully direct route -- in which the Galois factors themselves follow from
the $A_5/\mathbb{Z}_3$ structure -- is the remaining open edge.

\paragraph{What is genuinely new.}
Prior to this document, Doc.~293 had computed the $p$-values but had not
identified their ratios as exact $\phi$-powers. The identities
$p_1/p_2 = \phi^8$ and $p_0/p_2 = 2\phi^4$ are new. The observation
$m_\tau/m_e \approx 74\times(p_1/p_2)$ connects, for the first time,
the $A_5$ geometry of Doc.~293 to the Galois structure of Doc.~338 through
a concrete mass ratio.
